\documentclass[aspectratio=169,10pt]{beamer}
% Shared Beamer setup for Fall 2026 course slides.
\mode<presentation>{
  \usetheme{Madrid}
  \usecolortheme{default}
}

\definecolor{CalPolyGreen}{HTML}{154734}
\definecolor{CalPolyGold}{HTML}{BD8B13}
\definecolor{DeckInk}{HTML}{1F2933}
\definecolor{DeckMuted}{HTML}{5F6B7A}

\setbeamercolor{structure}{fg=CalPolyGreen}
\setbeamercolor{palette primary}{bg=CalPolyGreen,fg=white}
\setbeamercolor{palette secondary}{bg=DeckInk,fg=white}
\setbeamercolor{palette tertiary}{bg=CalPolyGold,fg=DeckInk}
\setbeamercolor{frametitle}{bg=CalPolyGreen,fg=white}
\setbeamercolor{title}{fg=CalPolyGreen}
\setbeamercolor{normal text}{fg=DeckInk,bg=white}
\setbeamercolor{block title}{bg=CalPolyGreen,fg=white}
\setbeamercolor{block body}{bg=black!3,fg=DeckInk}

\setbeamertemplate{navigation symbols}{}
\setbeamertemplate{footline}[frame number]
\setbeamertemplate{itemize item}{\color{CalPolyGold}$\blacktriangleright$}
\setbeamertemplate{itemize subitem}{\color{CalPolyGreen}$\bullet$}

\newcommand{\CourseNumber}{Course}
\newcommand{\CourseTitle}{Course Title}
\newcommand{\LectureNumber}{0}
\newcommand{\LectureTitle}{Lecture Title}
\newcommand{\LectureDate}{Date}
\newcommand{\CourseUnit}{Unit}

\newcommand{\coursenumber}[1]{\renewcommand{\CourseNumber}{#1}}
\newcommand{\coursetitle}[1]{\renewcommand{\CourseTitle}{#1}}
\newcommand{\lecturenumber}[1]{\renewcommand{\LectureNumber}{#1}}
\newcommand{\lecturetitle}[1]{\renewcommand{\LectureTitle}{#1}}
\newcommand{\lecturedate}[1]{\renewcommand{\LectureDate}{#1}}
\newcommand{\courseunit}[1]{\renewcommand{\CourseUnit}{#1}}

\author{Kirk Duran}
\institute{California Polytechnic State University, San Luis Obispo}
\date{\LectureDate}

\newcommand{\makedecktitle}{
  \begin{frame}[plain]
    \vspace{0.25in}
    {\usebeamerfont{title}\color{CalPolyGreen}\LectureTitle\par}
    \vspace{0.18in}
    {\large \CourseNumber{}: \CourseTitle\par}
    \vspace{0.08in}
    {\large Lecture \LectureNumber{} -- \LectureDate\par}
    \vfill
    {\small Kirk Duran \textbar{} Cal Poly, San Luis Obispo\par}
  \end{frame}
}

\newcommand{\placeholdernote}{
  \begin{block}{Placeholder}
    Replace these bullets with the final lecture material, examples, and in-class activities.
  \end{block}
}

\AtBeginSection[]{
  \begin{frame}{Roadmap}
    \tableofcontents[currentsection]
  \end{frame}
}

\graphicspath{{../assets/lecture-03/}}
\setbeamersize{text margin left=0.42in,text margin right=0.42in}

\coursenumber{CS 4880}
\coursetitle{Artificial Intelligence}
\lecturenumber{3}
\lecturetitle{Intelligent Agents II}
\lecturedate{August 28, 2026}
\courseunit{1. Intelligent Agents}

\begin{document}

\makedecktitle

% ============================================================
% Opening and Review
% ============================================================

\begin{frame}[shrink=10]{Artificial Intelligence And Human Cognition}
  \begin{keyidea}
    AI mirrors human abilities: perception, memory, planning, and learning.
  \end{keyidea}
  \vspace{0.12in}
  \begin{itemize}
    \item AI aims to reproduce useful forms of intelligence and adaptability in machines.
    \item Definitions, models, and techniques in AI often reflect how we think humans reason and learn.
    \item The goal is not necessarily to copy biology exactly; it is to build systems that exhibit useful intelligent behavior.
    \item You already have an intuitive sense of many concepts that we will now formalize.
  \end{itemize}
\end{frame}

\begin{frame}{Review: Intelligent Agents}
  \begin{cpdefinition}
    An \textbf{agent} perceives an environment through sensors and acts upon that environment through actuators.
  \end{cpdefinition}
  \vspace{0.12in}
  \begin{center}
    {\Large Environment $\rightarrow$ Percepts $\rightarrow$ Agent $\rightarrow$ Actions $\rightarrow$ Environment}
  \end{center}
  \vspace{0.12in}
  \begin{itemize}
    \item A percept is what the agent senses at one moment.
    \item A percept sequence is the complete history of what it has sensed.
    \item The \textbf{agent function} maps percept histories to actions: $f:P^*\rightarrow A$.
    \item A rational agent chooses actions expected to maximize its performance measure given the information available.
  \end{itemize}
\end{frame}

% ============================================================
\sectionframe{Part 1: From Agent Functions To Agent Programs}
% ============================================================

\begin{frame}[fragile]{The Table-Driven Approach}
  \begin{columns}[T,onlytextwidth]
    \begin{column}{0.48\textwidth}
      \begin{keyidea}
        Imagine a perfect lookup table containing the correct action for every possible percept sequence.
      \end{keyidea}
      \vspace{0.08in}
      \begin{itemize}
        \item Append the newest percept to the percept history.
        \item Use the complete percept sequence as a table key.
        \item Return the action stored in that entry.
        \item In principle, the agent needs no run-time reasoning or learning.
      \end{itemize}
    \end{column}
    \begin{column}{0.47\textwidth}
\begin{verbatim}
function TABLE-DRIVEN-AGENT(percept):
    persistent percepts = []
    persistent table

    append percept to percepts
    action = LOOKUP(percepts, table)
    return action
\end{verbatim}
    \end{column}
  \end{columns}
\end{frame}

\begin{frame}[shrink=35]{Reality Check: Why A Table Does Not Work}
  \begin{block}{Storage Explosion}
    The number of possible percept histories grows combinatorially or exponentially. Rich sensors make a realistic table astronomically large.
  \end{block}
  \begin{block}{Impossible To Construct}
    Someone would need to precompute the correct action for every situation the agent could ever encounter.
  \end{block}
  \begin{block}{No Generalization}
    A static table treats similar situations as unrelated entries instead of exploiting common structure.
  \end{block}
  \begin{block}{No Room For Learning}
    New situations require new table entries; the representation itself provides no mechanism for adaptation.
  \end{block}
  \vspace{0.06in}
  \begin{keyidea}
    We need compact programs that reason from structure rather than enumerate every possible history.
  \end{keyidea}
\end{frame}

\begin{frame}{Kinds Of Agent Programs}
  \begin{enumerate}
    \item \textbf{Simple reflex agents}: react to the current percept.
    \item \textbf{Model-based reflex agents}: maintain internal state.
    \item \textbf{Goal-based agents}: reason about desired future states.
    \item \textbf{Utility-based agents}: compare how desirable possible outcomes are.
    \item \textbf{Learning agents}: improve behavior from experience.
  \end{enumerate}
  \vspace{0.14in}
  \begin{keyidea}
    These are increasingly expressive design patterns, not mutually exclusive product categories. Real systems often combine them.
  \end{keyidea}
\end{frame}

% ============================================================
\sectionframe{Part 2: Simple Reflex Agents}
% ============================================================

\begin{frame}{Is This Intelligence?}
  \begin{columns}[T,onlytextwidth]
    \begin{column}{0.58\textwidth}
      \imageplaceholder{stove-reflex.jpg}
    \end{column}
    \begin{column}{0.36\textwidth}
      \begin{activity}
        A sensor detects a condition and immediately triggers a response. Is that intelligent behavior?
      \end{activity}
      \vspace{0.1in}
      \begin{itemize}
        \item It senses.
        \item It acts.
        \item It may satisfy an objective.
        \item It may not remember, plan, or learn.
      \end{itemize}
    \end{column}
  \end{columns}
\end{frame}

\begin{frame}{Is This Intelligent?}
  \begin{columns}[T,onlytextwidth]
    \begin{column}{0.42\textwidth}
      \imageplaceholder{water-fountain.jpg}
    \end{column}
    \begin{column}{0.52\textwidth}
      \begin{activity}
        An automatic bottle filler detects a bottle and turns on the water. What would make this an agent? What would make it intelligent?
      \end{activity}
      \vspace{0.12in}
      \begin{keyidea}
        Intelligence is not necessarily all-or-nothing. Useful agent behavior can emerge from extremely simple decision rules.
      \end{keyidea}
    \end{column}
  \end{columns}
\end{frame}

\begin{frame}{Simple Reflex Agent}
  \begin{cpdefinition}
    A \textbf{simple reflex agent} chooses an action using only the current percept.
  \end{cpdefinition}
  \vspace{0.08in}
  \begin{columns}[T,onlytextwidth]
    \begin{column}{0.48\textwidth}
      \begin{block}{Characteristics}
        \begin{itemize}
          \item Ignores percept history.
          \item Maintains no internal world state.
          \item Uses \textbf{condition--action rules}.
        \end{itemize}
      \end{block}
      \begin{exampleblockcp}
        \textbf{IF} obstacle ahead \textbf{THEN} turn right.
      \end{exampleblockcp}
    \end{column}
    \begin{column}{0.46\textwidth}
      \begin{block}{Flow}
        \begin{enumerate}
          \item Perceive the environment.
          \item Match the percept to a rule condition.
          \item Execute the associated action.
          \item Repeat with no memory of the previous decision.
        \end{enumerate}
      \end{block}
    \end{column}
  \end{columns}
\end{frame}

\begin{frame}{Simple Reflex Agents: Strengths And Weaknesses}
  \begin{columns}[T,onlytextwidth]
    \begin{column}{0.48\textwidth}
      \begin{block}{Advantages}
        \begin{itemize}
          \item Fast and computationally cheap.
          \item Requires little or no memory.
          \item Easy to implement and debug.
          \item Works well when the environment is fully observable and the needed information is available now.
        \end{itemize}
      \end{block}
    \end{column}
    \begin{column}{0.48\textwidth}
      \begin{block}{Downside}
        \begin{itemize}
          \item Missing information can cause incorrect actions.
          \item The same percept may require different actions in different hidden states.
          \item Local rules can create brittle or repetitive behavior.
        \end{itemize}
      \end{block}
    \end{column}
  \end{columns}
  \vspace{0.08in}
  \begin{center}
    \imageplaceholder{simple-reflex-diagram.jpg}
  \end{center}
\end{frame}

\begin{frame}[shrink=30]{Simple Reflex Example: Thermostat}
  \begin{columns}[T,onlytextwidth]
    \begin{column}{0.54\textwidth}
      \begin{block}{Percept}
        Current room temperature.
      \end{block}
      \begin{block}{Condition--Action Rules}
        \begin{itemize}
          \item If temperature is below the lower threshold, turn on heat.
          \item If temperature is within the desired band, maintain the current state.
          \item If temperature is above the upper threshold, turn off heat.
        \end{itemize}
      \end{block}
      \begin{keyidea}
        The action depends on the current temperature, not a detailed history of the room.
      \end{keyidea}
    \end{column}
    \begin{column}{0.4\textwidth}
      \imageplaceholder{thermostat-visual.jpg}
    \end{column}
  \end{columns}
\end{frame}

\begin{frame}[shrink=35]{Historical Reflex Example: Jacquard Loom}
  \begin{columns}[T,onlytextwidth]
    \begin{column}{0.55\textwidth}
      \begin{itemize}
        \item Joseph-Marie Jacquard's early nineteenth-century loom used punched cards to control weaving patterns.
        \item A hole or no-hole pattern directly determined which threads were lifted.
        \item The loom did not represent the textile's purpose or ``understand'' the pattern.
        \item Its behavior was a direct response to the current card pattern.
      \end{itemize}
      \vspace{0.1in}
      \begin{exampleblockcp}
        Hole present $\rightarrow$ lift thread.\newline
        No hole $\rightarrow$ do nothing.
      \end{exampleblockcp}
      \vspace{0.06in}
      \begin{keyidea}
        A useful system can implement precise conditional behavior without a rich internal model.
      \end{keyidea}
    \end{column}
    \begin{column}{0.39\textwidth}
      \imageplaceholder{jacquard-loom.jpg}
    \end{column}
  \end{columns}
\end{frame}

% ============================================================
\sectionframe{Part 3: Model-Based Reflex Agents}
% ============================================================

\begin{frame}[shrink=25]{Model-Based Reflex Agent}
  \begin{cpdefinition}
    A \textbf{model-based reflex agent} maintains an internal state so it can reason about information that is not directly observable now.
  \end{cpdefinition}
  \vspace{0.08in}
  \begin{columns}[T,onlytextwidth]
    \begin{column}{0.47\textwidth}
      \begin{block}{Transition Model}
        Describes how the world changes over time and how the agent's actions change it.
      \end{block}
      \begin{exampleblockcp}
        Turning the steering wheel changes a vehicle's heading.
      \end{exampleblockcp}
    \end{column}
    \begin{column}{0.47\textwidth}
      \begin{block}{Sensor Model}
        Describes how states of the world produce percepts.
      \end{block}
      \begin{exampleblockcp}
        Rain on a camera lens can distort the observed image.
      \end{exampleblockcp}
    \end{column}
  \end{columns}
  \vspace{0.08in}
  \begin{keyidea}
    Previous state + previous action + current percept $\rightarrow$ inferred current state $\rightarrow$ condition--action rule.
  \end{keyidea}
\end{frame}

\begin{frame}{Model-Based Reflex Example: Self-Driving Car}
  \begin{columns}[T,onlytextwidth]
    \begin{column}{0.48\textwidth}
      \begin{itemize}
        \item Track nearby vehicles even when they are temporarily occluded.
        \item Predict how those vehicles are likely to move.
        \item Estimate position, speed, lane, and surrounding road geometry.
        \item Base decisions on maintained state rather than one sensor frame alone.
      \end{itemize}
      \vspace{0.08in}
      \begin{block}{Downside}
        More memory and computation are required, and the agent can still make mistakes when its model is inaccurate or incomplete.
      \end{block}
    \end{column}
    \begin{column}{0.46\textwidth}
      \imageplaceholder{model-based-diagram.jpg}
    \end{column}
  \end{columns}
\end{frame}

\begin{frame}{Reflex Or Model-Based Reflex?}
  \begin{columns}[T,onlytextwidth]
    \begin{column}{0.58\textwidth}
      \imageplaceholder{vacuum-reflex-or-model.jpg}
    \end{column}
    \begin{column}{0.36\textwidth}
      \begin{activity}
        A robot vacuum maps rooms, remembers cleaned areas, and tracks its location. Is it still a simple reflex agent?
      \end{activity}
      \vspace{0.1in}
      Look for:
      \begin{itemize}
        \item an internal map,
        \item memory of earlier percepts,
        \item prediction of hidden state,
        \item rules operating on inferred state.
      \end{itemize}
    \end{column}
  \end{columns}
\end{frame}

% ============================================================
\sectionframe{Part 4: Goal-Based And Utility-Based Agents}
% ============================================================

\begin{frame}{Goal-Based Agent}
  \begin{cpdefinition}
    A \textbf{goal-based agent} evaluates actions by whether their consequences move the world toward a desired goal state.
  \end{cpdefinition}
  \vspace{0.08in}
  \begin{columns}[T,onlytextwidth]
    \begin{column}{0.5\textwidth}
      \begin{itemize}
        \item Uses a model of how actions change the world.
        \item Adds explicit information about which outcomes are desirable.
        \item Can reason about the future instead of reacting immediately.
        \item Often uses search or planning to find a sequence of actions.
      \end{itemize}
    \end{column}
    \begin{column}{0.44\textwidth}
      \begin{block}{Flow}
        Perceive $\rightarrow$ update state $\rightarrow$ consult goal $\rightarrow$ search / plan $\rightarrow$ act $\rightarrow$ re-plan as needed.
      \end{block}
    \end{column}
  \end{columns}
\end{frame}

\begin{frame}{Goal-Based Architecture}
  \begin{columns}[T,onlytextwidth]
    \begin{column}{0.55\textwidth}
      \imageplaceholder{goal-based-diagram.jpg}
    \end{column}
    \begin{column}{0.39\textwidth}
      \begin{itemize}
        \item Goals describe what success looks like.
        \item The model predicts the consequences of actions.
        \item Search or planning finds a path from the current state to a goal state.
      \end{itemize}
      \vspace{0.1in}
      \begin{keyidea}
        A goal-based agent can choose actions that are useful only because of where they lead later.
      \end{keyidea}
    \end{column}
  \end{columns}
\end{frame}

\begin{frame}{Goal-Based Example: Vacuum Robot}
  \begin{columns}[T,onlytextwidth]
    \begin{column}{0.45\textwidth}
      \imageplaceholder{goal-vacuum.jpg}
    \end{column}
    \begin{column}{0.5\textwidth}
      \begin{itemize}
        \item Perceive dirt, obstacles, and current position.
        \item Update an internal map from previous movements and percepts.
        \item Check whether the goal is achieved: is the required area clean?
        \item If not, select an action or path that moves closer to the goal.
        \item Execute and repeat until the goal is met.
      \end{itemize}
    \end{column}
  \end{columns}
\end{frame}

\begin{frame}[shrink=10]{Is A Goal Enough?}
  \begin{columns}[T,onlytextwidth]
    \begin{column}{0.42\textwidth}
      \imageplaceholder{chess-goal.jpg}
    \end{column}
    \begin{column}{0.52\textwidth}
      \begin{activity}
        A chess agent wants to win. Is ``win the game'' enough information to choose the next move?
      \end{activity}
      \vspace{0.08in}
      \begin{itemize}
        \item Win quickly or slowly?
        \item Keep more valuable pieces?
        \item Avoid risk?
        \item Prefer one winning line over another?
      \end{itemize}
      \vspace{0.08in}
      \begin{keyidea}
        Goals often give a binary answer: acceptable or not. They do not necessarily rank alternatives.
      \end{keyidea}
    \end{column}
  \end{columns}
\end{frame}

\begin{frame}{Utility-Based Example: Chess}
  \begin{keyidea}
    Even when the ultimate goal is ``win,'' an agent can evaluate intermediate states by estimating how desirable they are.
  \end{keyidea}
  \vspace{0.1in}
  Possible evaluation features include:
  \begin{columns}[T,onlytextwidth]
    \begin{column}{0.47\textwidth}
      \begin{itemize}
        \item preserve valuable pieces,
        \item control the center,
        \item protect the king,
      \end{itemize}
    \end{column}
    \begin{column}{0.47\textwidth}
      \begin{itemize}
        \item avoid weak pawn structures,
        \item maintain mobility,
        \item increase the probability of eventually winning.
      \end{itemize}
    \end{column}
  \end{columns}
  \vspace{0.1in}
  \begin{exampleblockcp}
    A utility value compresses several strategic preferences into a quantity the agent can compare.
  \end{exampleblockcp}
\end{frame}

\begin{frame}{Utility-Based Agent}
  \begin{cpdefinition}
    A \textbf{utility function} assigns a numerical value to states or outcomes representing how desirable they are to the agent.
  \end{cpdefinition}
  \vspace{0.08in}
  \begin{columns}[T,onlytextwidth]
    \begin{column}{0.48\textwidth}
      \begin{itemize}
        \item Ranks alternatives instead of treating success as yes/no.
        \item Handles conflicting objectives.
        \item Under uncertainty, compares actions by \textbf{expected utility}.
      \end{itemize}
      \vspace{0.08in}
      \begin{center}
        $EU(a)=\sum_s P(s\mid a,e)U(s)$
      \end{center}
    \end{column}
    \begin{column}{0.46\textwidth}
      \imageplaceholder{utility-based-diagram.jpg}
    \end{column}
  \end{columns}
\end{frame}

\begin{frame}[shrink=10]{Utility-Based Example: Self-Driving Taxi}
  A rational taxi may optimize several objectives simultaneously:
  \begin{itemize}
    \item shortest travel time,
    \item least fuel or energy consumption,
    \item lowest accident risk,
    \item passenger comfort,
    \item compliance with traffic laws.
  \end{itemize}
  \vspace{0.1in}
  \begin{block}{Why Utility Helps}
    It supports tradeoffs such as speed versus safety and allows the agent to compare several acceptable outcomes.
  \end{block}
  \begin{block}{Tradeoff}
    Utility-based reasoning can be computationally expensive because the agent may need to evaluate many possibilities and uncertain outcomes.
  \end{block}
\end{frame}

\begin{frame}[shrink=20]{Goals Or Utility Are Not Enough}
  \begin{columns}[T,onlytextwidth]
    \begin{column}{0.56\textwidth}
      \imageplaceholder{deepblue-alphago.jpg}
    \end{column}
    \begin{column}{0.38\textwidth}
      \begin{itemize}
        \item Deep Blue used massive search plus carefully engineered evaluation methods to defeat Garry Kasparov.
        \item AlphaGo combined search with learned policy and value models to defeat Lee Sedol in 2016.
        \item A goal or utility function tells an agent what it wants; it does not automatically tell the agent how to act well.
      \end{itemize}
    \end{column}
  \end{columns}
  \vspace{0.08in}
  \begin{keyidea}
    Intelligent systems often combine representation, models, search, utility, and learning.
  \end{keyidea}
\end{frame}

% ============================================================
\sectionframe{Part 5: Learning Agents}
% ============================================================

\begin{frame}{Learning Agents}
  \begin{cpdefinition}
    A \textbf{learning agent} improves its behavior over time based on experience or feedback.
  \end{cpdefinition}
  \vspace{0.08in}
  \begin{enumerate}
    \item Perceive the environment.
    \item Take an action using the current agent program.
    \item Observe the outcome.
    \item Use feedback to improve internal models, rules, utilities, or strategies.
    \item Repeat with better knowledge next time.
  \end{enumerate}
  \vspace{0.08in}
  \begin{keyidea}
    Learning allows an agent to begin with incomplete knowledge, adapt to an unfamiliar environment, and improve through experience.
  \end{keyidea}
\end{frame}

\begin{frame}{Four Components Of A Learning Agent}
  \begin{columns}[T,onlytextwidth]
    \begin{column}{0.23\textwidth}
      \begin{block}{Performance Element}
        Chooses the actions that interact with the external world.
      \end{block}
    \end{column}
    \begin{column}{0.23\textwidth}
      \begin{block}{Learning Element}
        Modifies the performance element to improve behavior from experience.
      \end{block}
    \end{column}
    \begin{column}{0.23\textwidth}
      \begin{block}{Critic}
        Evaluates how well the agent is performing relative to the performance measure.
      \end{block}
    \end{column}
    \begin{column}{0.23\textwidth}
      \begin{block}{Problem Generator}
        Suggests exploratory actions that may produce informative experience.
      \end{block}
    \end{column}
  \end{columns}
  \vspace{0.1in}
  \begin{keyidea}
    Learning separates acting, evaluating, improving, and exploring.
  \end{keyidea}
\end{frame}

\begin{frame}{Learning Agent Example: Self-Driving Car}
  \begin{itemize}
    \item \textbf{Performance element:} chooses braking, acceleration, and steering actions.
    \item \textbf{Critic:} evaluates whether braking was too hard, too soft, unsafe, or uncomfortable.
    \item \textbf{Learning element:} updates a transition model, control policy, or prediction model based on feedback.
    \item \textbf{Problem generator:} encourages informative experience, such as varied braking maneuvers under safe conditions.
  \end{itemize}
  \vspace{0.1in}
  \begin{exampleblockcp}
    Observed outcome $\rightarrow$ feedback $\rightarrow$ model or policy update $\rightarrow$ improved future action.
  \end{exampleblockcp}
\end{frame}

% ============================================================
\sectionframe{Part 6: Problem Formulation And Abstraction}
% ============================================================

\begin{frame}[shrink=20]{Problem Abstraction}
  \begin{columns}[T,onlytextwidth]
    \begin{column}{0.51\textwidth}
      \begin{cpdefinition}
        \textbf{Abstraction} represents a real problem using only the details needed for reasoning about a solution.
      \end{cpdefinition}
      \vspace{0.08in}
      \begin{itemize}
        \item Real road travel includes weather, traffic, passengers, music, laws, road paint, and countless other facts.
        \item A navigation model may keep only locations, roads, travel costs, and constraints.
        \item Cities or intersections become \textbf{states}.
        \item Roads become \textbf{actions or transitions}.
        \item Irrelevant details are filtered out.
      \end{itemize}
    \end{column}
    \begin{column}{0.43\textwidth}
      \imageplaceholder{problem-abstraction-navigation.jpg}
    \end{column}
  \end{columns}
\end{frame}

\begin{frame}{Formal Search Problem Formulation}
  A search problem is commonly described by:
  \begin{itemize}
    \item \textbf{Initial state}: where the agent starts.
    \item \textbf{Actions}: the choices available in a state.
    \item \textbf{Transition model}: the state produced by taking an action.
    \item \textbf{Goal test}: whether a state satisfies the objective.
    \item \textbf{Path cost}: the accumulated cost of a sequence of actions.
  \end{itemize}
  \vspace{0.12in}
  \begin{keyidea}
    Search algorithms operate on this abstract state space rather than on every detail of the physical world.
  \end{keyidea}
\end{frame}

% ============================================================
\sectionframe{Part 7: Problem Types}
% ============================================================

\begin{frame}{Single-State Problem}
  \begin{cpdefinition}
    A \textbf{single-state problem} assumes the agent knows exactly which state it is in and can predict the result of each action.
  \end{cpdefinition}
  \vspace{0.12in}
  Typical assumptions:
  \begin{itemize}
    \item the environment is fully observable,
    \item the transition model is known,
    \item action outcomes are deterministic.
  \end{itemize}
  \vspace{0.12in}
  \begin{keyidea}
    The agent can plan a complete sequence of actions without uncertainty about its current state or the immediate effect of an action.
  \end{keyidea}
\end{frame}

\begin{frame}{Chess As A Single-State Problem}
  \begin{columns}[T,onlytextwidth]
    \begin{column}{0.43\textwidth}
      \imageplaceholder{chess-board.jpg}
    \end{column}
    \begin{column}{0.51\textwidth}
      \begin{itemize}
        \item The current board position is fully visible.
        \item Legal moves are known.
        \item Executing a legal move has a deterministic result.
        \item The agent therefore knows exactly what board state each move will produce.
      \end{itemize}
      \vspace{0.1in}
      \begin{keyidea}
        Chess is still strategically difficult because another agent chooses future moves; ``single-state'' refers to state certainty, not overall simplicity.
      \end{keyidea}
    \end{column}
  \end{columns}
\end{frame}

\begin{frame}{Multiple-State Problem}
  \begin{cpdefinition}
    In a \textbf{multiple-state problem}, the agent cannot identify one exact current world state from its percepts.
  \end{cpdefinition}
  \vspace{0.12in}
  \begin{itemize}
    \item The world may be partially observable.
    \item Several physical states may be consistent with the same percept history.
    \item The agent therefore reasons about a set of possible current states.
  \end{itemize}
  \vspace{0.12in}
  \begin{exampleblockcp}
    A robot with limited sensors may know that it is in one of several rooms without knowing exactly which room.
  \end{exampleblockcp}
\end{frame}

\begin{frame}[shrink=30]{Belief States}
  \begin{columns}[T,onlytextwidth]
    \begin{column}{0.48\textwidth}
      \imageplaceholder{multiple-state-map.jpg}
    \end{column}
    \begin{column}{0.46\textwidth}
      \begin{cpdefinition}
        A \textbf{belief state} represents all states the agent currently considers possible.
      \end{cpdefinition}
      \vspace{0.08in}
      \begin{itemize}
        \item New percepts eliminate inconsistent possibilities.
        \item Actions transform the set of possible states.
        \item Planning occurs over belief states when the true state is hidden.
      \end{itemize}
      \vspace{0.08in}
      \begin{keyidea}
        The agent reasons about what it \emph{believes} the world could be, not only one assumed world state.
      \end{keyidea}
    \end{column}
  \end{columns}
\end{frame}

\begin{frame}{Contingency Problems}
  \begin{cpdefinition}
    A \textbf{contingency problem} requires a plan that can react to different observations or outcomes encountered during execution.
  \end{cpdefinition}
  \vspace{0.08in}
  \begin{itemize}
    \item The environment may be nondeterministic: the same action can lead to different results.
    \item External agents or physical uncertainty may interfere with the expected result.
    \item The agent may not have complete control over the environment.
    \item New observations can force the plan to branch.
  \end{itemize}
  \vspace{0.08in}
  \begin{exampleblockcp}
    Drive toward the destination. If the highway is clear, continue. If an accident blocks it, use the alternate route. If both are blocked, stop and re-plan.
  \end{exampleblockcp}
\end{frame}

\begin{frame}{Exploration Problems}
  \begin{cpdefinition}
    In an \textbf{exploration problem}, the agent does not initially know enough about the environment to plan completely in advance.
  \end{cpdefinition}
  \vspace{0.08in}
  \begin{itemize}
    \item The agent may not know the full state space, available actions, or transition model.
    \item It gathers information by interacting with the environment.
    \item It must build or refine a model while simultaneously attempting to achieve its goal.
    \item Exploration trades immediate performance for information that may improve future decisions.
  \end{itemize}
  \vspace{0.08in}
  \begin{keyidea}
    The agent learns how the world works by acting in it.
  \end{keyidea}
\end{frame}

\begin{frame}[shrink=10]{Exploration Example: Self-Driving Car}
  \begin{columns}[T,onlytextwidth]
    \begin{column}{0.48\textwidth}
      \begin{itemize}
        \item The car enters an unmapped city or road network.
        \item It does not begin with complete knowledge of the roads.
        \item It explores streets and intersections while localizing itself.
        \item It incrementally builds a map and improves its routing.
      \end{itemize}
      \vspace{0.1in}
      \begin{exampleblockcp}
        Simultaneous localization and mapping (SLAM) is a classic example of building a model while navigating.
      \end{exampleblockcp}
    \end{column}
    \begin{column}{0.46\textwidth}
      \imageplaceholder{exploration-car.jpg}
    \end{column}
  \end{columns}
\end{frame}

\begin{frame}{Problem Domains}
  \begin{columns}[T,onlytextwidth]
    \begin{column}{0.47\textwidth}
      \begin{block}{Puzzles And Mazes}
        \begin{itemize}
          \item Sliding-tile puzzle such as the 8-puzzle.
          \item A maze with full observability.
          \item A maze where the agent does not know its starting location.
          \item A partially completed jigsaw puzzle without edge pieces.
        \end{itemize}
      \end{block}
    \end{column}
    \begin{column}{0.47\textwidth}
      \begin{block}{Strategy And Simulation}
        \begin{itemize}
          \item Minesweeper.
          \item Battleship.
          \item StarCraft.
          \item Escape rooms and other interactive environments.
        \end{itemize}
      \end{block}
    \end{column}
  \end{columns}
  \vspace{0.08in}
  \begin{activity}
    For each problem, identify what the agent knows, what it does not know, and whether outcomes are deterministic.
  \end{activity}
\end{frame}

% ============================================================
\sectionframe{Part 8: Putting The Architectures Together}
% ============================================================

\begin{frame}{Agent Architectures: What Changes?}
  \scriptsize
  \begin{center}
  \begin{tabular}{@{}p{0.19\linewidth} p{0.2\linewidth} p{0.25\linewidth} p{0.24\linewidth}@{}}
    \textbf{Architecture} & \textbf{Uses} & \textbf{Main Question} & \textbf{Key Addition} \\
    \hline
    Simple reflex & Current percept & What should I do now? & Condition--action rules \\
    Model-based reflex & Percept + internal state & What is the world like now? & Memory and model \\
    Goal-based & State + goals & How can I reach success? & Search / planning \\
    Utility-based & State + utilities & Which outcome is best? & Preferences and tradeoffs \\
    Learning & Experience + feedback & How can I improve? & Adaptation over time \\
  \end{tabular}
  \end{center}
\end{frame}

\begin{frame}[shrink=10]{A Modern Agent Is Usually Hybrid}
  \begin{keyidea}
    Real AI systems commonly combine several of today's architectures rather than choosing exactly one.
  \end{keyidea}
  \vspace{0.12in}
  A self-driving car may:
  \begin{itemize}
    \item use reflexes for emergency braking,
    \item maintain a model of surrounding traffic,
    \item plan toward a destination,
    \item optimize safety, time, energy, and comfort,
    \item learn from data and experience.
  \end{itemize}
  \vspace{0.12in}
  \begin{activity}
    Given an environment and a performance measure, what internal structure does the agent actually need?
  \end{activity}
\end{frame}

\begin{frame}{Next Time: Search Algorithms}
  \begin{columns}[T,onlytextwidth]
    \begin{column}{0.46\textwidth}
      \imageplaceholder{search-maze.jpg}
    \end{column}
    \begin{column}{0.48\textwidth}
      \begin{keyidea}
        Goal-based agents need a systematic way to find sequences of actions that move from an initial state to a goal state.
      \end{keyidea}
      \vspace{0.12in}
      \begin{itemize}
        \item State spaces.
        \item Search trees.
        \item Breadth-first and depth-first search.
        \item Cost and optimality.
        \item Heuristics later.
      \end{itemize}
    \end{column}
  \end{columns}
\end{frame}

\end{document}
